\subsection{Locally integrable functions}

PROPOSITION 1. --- Let g be a function defined locally almost everywhere in T (for the positive measure \(\mu\) ), with values in a Banach space F (resp. in \(\overline{R}\) ). The following properties are equivalent:

\begin{enumerate}
\def\labelenumi{\alph{enumi})}
\item
  For every point \(t \in \mathbb{T}\) , there exists a neighborhood \(\mathbb{V}\) of \(t\) such that the function \(\mathbf{g}\varphi_{\mathbb{V}}\) is \(\mu\) -integrable.
\item
  The function \(\mathbf{g}\) is \(\mu\) -measurable and, for every compact set \(\mathrm{K} \subset \mathrm{T}\) , \(\int^{*} |\mathbf{g}| \varphi_{\mathrm{K}} d\mu < +\infty\) .
\item
  For every numerical function \(h \in \mathcal{K}(\mathrm{T})\) , gh is \(\mu\) -integrable.
\end{enumerate}

Let us show that a) implies b); for, the function g is measurable by the principle of localization (Ch. IV, §5, No.~2, Prop. 4). On the other hand, for every \(t \in \mathbf{K}\) there exists, by hypothesis, a neighborhood \(V_t\) of \(t\) in \(T\) such that \(\mathbf{g}\varphi_{V_t}\) is integrable; one can therefore cover \(K\) by a finite number of neighborhoods \(V_i\) ( \(1 \leqslant i \leqslant n\) ) such that the functions \(\mathbf{g}\varphi_{V_i}\) are integrable. Since \(|\mathbf{g}| \varphi_K \leqslant \sum_{i=1}^{n} |\mathbf{g}| \varphi_{V_i}\) , one has \(\int^* |\mathbf{g}| \varphi_K d\mu < +\infty\) .

Secondly, b) implies c), because gh is then measurable, and if L is the compact support of h then \(|gh| \leqslant \|h\|\cdot|g|\varphi_{L}\) , therefore \(\int^{*}|gh|d\mu < +\infty\) by hypothesis; gh is therefore integrable by the criterion for integrability (Ch. IV, §5, No.~6, Th. 5).

Finally, c) implies a). Indeed, for every \(t \in T\) let V be a compact neighborhood of t. There exists a continuous mapping h of T into {[}0,1{]}, equal to 1 on V and with compact support (Ch. III, §1, No.~2, Lemma 1); by hypothesis gh is integrable, therefore so is \(\mathbf{g}\varphi_{\mathrm{V}} = (\mathbf{g}h)\varphi_{\mathrm{V}}\) (Ch. IV, §5, No.~6, Cor. 3 of Th. 5).

DEFINITION 1. --- A function g, defined locally almost everywhere in T (for the positive measure \(\mu\) ), with values in a Banach space F (resp. in \(\overline{R}\) ), is said to be locally integrable for \(\mu\) (or locally \(\mu\) -integrable) if it satisfies the conditions a), b), c) of Prop. 1. If \(\theta\) is a complex measure, a function g defined locally \(\theta\) -almost everywhere is said to be locally \(\theta\) -integrable if it is locally integrable for the positive measure \(|\theta|\) .

If \(\mathbf{g}\) is locally \(\theta\) -integrable, then every function equal to \(\mathbf{g}\) locally almost everywhere is locally integrable. It is clear that the sum of two locally integrable functions is locally integrable. The functions with values in \(\mathbf{F}\) , everywhere defined and locally integrable for \(\theta\) , form a vector space denoted \(\mathcal{L}_{\mathrm{loc}}^{1}(\mathrm{T},\theta ;\mathrm{F})\) ; when \(\mathbf{F} = \mathbf{R}\) or \(\mathbf{C}\) , the mention of \(\mathbf{F}\) is often omitted if there is no ambiguity. This space will always be equipped (absent express mention to the contrary) with the topology defined by the semi-norms \(\mathbf{g}\mapsto \int |\mathbf{g}\varphi_{\mathbf{K}}|d|\theta |\) , where \(\mathbf{K}\) runs over the set of compact subsets of \(\mathbf{T}\) . The associated Hausdorff space, the quotient of \(\mathcal{L}_{\mathrm{loc}}^{1}(\mathrm{T},\theta ;\mathrm{F})\) by the subspace \(\mathcal{N}_{\mathrm{F}}^{\infty}\) of mappings that are zero locally almost everywhere, is denoted \(\mathrm{L}_{\mathrm{loc}}^{1}(\mathrm{T},\theta ;\mathrm{F})\) . The spaces \(\mathrm{L}_{\mathrm{loc}}^{1}(\mathrm{T},\theta ;\mathrm{F})\) and \(\mathrm{L}_{\mathrm{loc}}^{1}(\mathrm{T},|\theta |; \mathrm{F})\) are identical.

It can be shown that the topological vector spaces just defined are complete (Exer. 31).

Every measurable function g, that is essentially bounded on every compact set, is locally integrable. For every number p such that \(1 \leqslant p \leqslant +\infty\) , every function \(g \in L_{F}^{p}\) is locally integrable; indeed, for every function \(h \in \mathcal{K}(T)\) , h belongs to \(L^{q}\) (where q is the exponent conjugate to p), therefore gh is integrable (Ch. IV, §6, No.~4, Cor. 4 of Th. 2).

Let \(F, G, H\) be three Banach spaces, and \((\mathbf{u}, \mathbf{v}) \mapsto \Phi(\mathbf{u}, \mathbf{v})\) a continuous bilinear mapping of \(F \times G\) into \(H\) . If \(f\) is locally integrable and takes its values in F, and if \(g \in L_{G}^{\infty}\) , then \(\Phi(\mathbf{f}, \mathbf{g})\) is locally integrable (Ch. IV, §6, No.~4, Cor. 1 of Th. 2).

